% !PRINTED 397
\weberpage{397}
Daß diese Zahl wirklich die genaue Anzahl der Geschlechter ist, werden wir später beweisen. Wir bemerken hierzu noch folgendes:

Die Tatsache, daß es für jedes $\Delta$ Zahlen $a$ gibt, für die $(\Delta, a) = -1$ ist, folgt aus dem Reziprozitätsgesetze der quadratischen Reste. Kann man, ohne dieses Gesetz vorauszusetzen, die Existenz solcher Zahlen $a$ nachweisen, so läßt sich umgekehrt dadurch ein Beweis des Reziprozitätsgesetzes ableiten. In dieser Weise hat \gesperrt{Gauss} seinen zweiten Beweis dieses Gesetzes hergeleitet. (Vgl. \gesperrt{Dirichlet-Dedekind}, \S\,152\,ff.)

Ist $r$ ein Normenrest und $\delta$ irgend ein Stammteiler von $D$, so ist
\weq{6}{N(\omega) = \frac{x^2 - y^2 D}{4} \equiv r \ (\mathrm{mod}\ \delta)}
und folglich, wenn $r$ relativ prim zu $\delta$ ist,
\weq{7}{(\delta, r) = +1.}

Für ungerade $\delta$ ist dies evident, für die geraden folgt es leicht in den verschiedenen Fällen von \S\,108. Wir setzen also allgemein, wenn $z$ relativ prim zu $D$ ist:
\weq{8}{\chi(z) = (\delta, z).}

Dann ist $(\delta, r) = 1$ für jede Zahl $r$ der Gruppe der absoluten Normenreste, und allgemein ist
\weq{9}{\chi(z) \chi(z_1) = \chi(z z_1).}

Für alle Ideale $\mathfrak{a}$ des \gesperrt{Hauptgeschlechtes}, die zu $D$ teilerfremd sind, ist
\[
\chi[N(\mathfrak{a})] = +1,
\]
und daraus folgt, daß $\chi[N(\mathfrak{a})]$ \gesperrt{für alle Ideale $\mathfrak{a}$ eines Geschlechtes einen und denselben Wert $\chi(G)$} hat.

\gesperrt{Diese Funktionen sind also die Charaktere der Gruppe der Geschlechter.}

\begin{center}\rule{0.32\linewidth}{0.4pt}\end{center}
